\section{Further uniform consequences of the arithmetic estimates}
\label{p1:uniform:appendix}

The following consequences record uniformities useful when the
counting argument is applied to selected subfamilies or to
conditional prime-sign systems. Their parameters and the
quantities being summed are kept explicit.

\subsection{Full-support relations at all cardinalities}

For \(r\ge2\), let \(\mathcal H_{r,N}\) be the ordered
lists of distinct, mutually separated regular windows for
which
\[
 s_i\in S_{x_i}\setminus\{0\},\qquad
                     s_1+\cdots+s_r=0
\]
has a solution. Minimality and affine compatibility are
not required.

\begin{proposition}[A weighted count of every full support]
\label{p1:uniform:full-support}
For fixed \(s\ge2\), \(u>0\), and \(\varepsilon>0\),
\[
 \sum_{r=s}^{|G|}\frac{u^r}{r!}
     \sum_{\boldsymbol x\in\mathcal H_{r,N}}p^r
       \ll_{s,u,C,\varepsilon}N^{-2s/(2s+1)+\varepsilon}.
\]
\end{proposition}

\begin{proof}
A full-support relation selects at least two true
values in every window, since selections are even
and nonzero. Consequently the moderate certificates
and the saturation graph of
Section~\ref{p14:cert:section} apply to these families,
without any primitivity assumption.

Use the parameters in
\eqref{p1:core:primitive-parameters}.
The total independent weight of moderate
certificates is \(N^{-2\beta+4/A_0+o(1)}
=N^{-91/50+o(1)}\), by
Lemma~\ref{p1:core:moderate}.
If a family has no moderate certificate, every
unmarked vertex has a neighbour in the saturation
graph. The tilted generating function of
Lemma~\ref{p1:core:generator}, now with no factor
for rank loss, bounds its weight beyond size \(s_0\)
by \(O_u(N^{-s_0\delta/2})\).
Choose a fixed \(s_0>\max(s,5000)\).
Both large-size contributions are smaller than
\(N^{-1+\varepsilon}\).

For each of the finitely many \(s\le r<s_0\),
Proposition~\ref{p1:arith:fixed-support} gives
\[
 |\mathcal H_{r,N}|
       \ll_{r,C,\varepsilon}
             N^{r-2r/(2r+1)+\varepsilon}.
 \]
Multiplying by \((up)^r/r!\), and using
\(2r/(2r+1)\ge2s/(2s+1)\), completes the proof.
The factor \(1/r!\) counts unordered sets; it is
not a bound for the growth of an arithmetic
constant in \(r\).
\end{proof}

\begin{lemma}[When no relation uses every block]
\label{p1:uniform:missing-block}
If a family \(I\) of \(r\) regular windows has no
relative relation meeting every block, then
\[
 q(I)\le p\sum_{x\in I}q(I\setminus\{x\}).
\]
Summed over any subfamily of ordered distinct
\(r\)-tuples, the resulting upper bound is at most
\(rNp\,\E(Z_G)_{r-1}\).
\end{lemma}

\begin{proof}
The relation space is covered by the kernels of
its \(r\) coordinate-block projections.
Hence \(2^{\rho(I)}\le\sum_x2^{\rho(I\setminus\{x\})}\).
If the full affine event is compatible, every
subevent is compatible and multiplication by
\(p^r\) proves the claim. If it is incompatible,
the left side is zero. Summing over the omitted
centre gives at most \(N\) choices, for each of
the \(r\) omitted positions.
\end{proof}

An elementary companion to such support bounds is the
following weighted triangle inequality. If
\((w_{xy})\) is a finite symmetric matrix of
nonnegative integers and \(W=\sum_{x,y}w_{xy}\), then
\[
 \sum_{x,y,z}\min(w_{xy},w_{xz},w_{yz})\le W^{3/2}.
\]
Indeed, for each integer \(h\ge1\), let \(A_h\) be the
symmetric zero-one matrix of entries at least \(h\),
and \(M_h\) its number of ones. Its triangle sum is
\(\operatorname{tr}(A_h^3)\le
(\operatorname{tr}(A_h^2))^{3/2}=M_h^{3/2}\).
Sum in \(h\) and use
\(\sum_hM_h^{3/2}\le(\sum_hM_h)^{3/2}\).

\subsection{Old square classes and a growing number of sources}

Let \(\mathcal P\) be a finite set of primes, and
\[
 A(X;\mathcal P)
   =\#\{n\le X:\sf(n)\text{ is supported on }\mathcal P\}.
\]

\begin{lemma}[Rankin under a logarithmic prime budget]
\label{p1:uniform:prime-budget}
For \(1/2<\sigma<1\),
\[
 A(X;\mathcal P)\le
 X^\sigma\zeta(2\sigma)
                  \prod_{\ell\in\mathcal P}(1+\ell^{-\sigma}).
\]
Uniformly for \(3/5\le\sigma\le3/4\), \(Y\ge3\), and
\(\sum_{\ell\in\mathcal P}\log\ell\le Y\),
\[
 \log A(X;\mathcal P)
       \le \sigma\log X+
              O\left(1+\frac{Y^{1-\sigma}}{\log Y}\right).
\]
\end{lemma}

\begin{proof}
Write \(n=du^2\), where \(d\) is squarefree and
supported on \(\mathcal P\), and sum
\((X/(du^2))^\sigma\).
For the second assertion, the prime-counting
upper bound and partial summation give
\(\sum_{\ell\le Y}\ell^{-\sigma}
\ll Y^{1-\sigma}/\log Y\) uniformly in the stated
compact interval. There are at most \(Y/\log Y\)
members of \(\mathcal P\) exceeding \(Y\);
their contribution satisfies the same bound.
Finally \(\zeta(2\sigma)\) is uniformly bounded.
\end{proof}

\begin{proposition}[Independent extensions at the critical order]
\label{p1:uniform:branching}
Put \(H=\log N\), \(h=\log H\), and fix \(K>0\).
For every \(j\le KH/h\) source windows, there is
an exceptional set of new centres of size at most
\(D_N^*\), where
\begin{equation}
 \log D_N^*
 \le\frac23H-\frac{H\log h}{9h}+O_{K,C}(H/h),
 \label{p1:uniform:branching-rate}
\end{equation}
such that outside it the \(L\) new relative bits
are exactly uniform and independent of all
previously exposed values.
\end{proposition}

\begin{proof}
Include in \(\mathcal P\) all primes at most \(B\)
and all primes dividing one of the \(jB\) old
values. Then
\[
 \sum_{\ell\in\mathcal P}\log\ell
       \ll_C B+jB\log(3N)\ll_{K,C}H^3/h.
\]
Take \(Y=C_{K,C}H^3/h\) and
\(\sigma=1-h/\log Y\). Eventually
\(\sigma\in[3/5,3/4]\) and \(Y^{1-\sigma}=H\).
Lemma~\ref{p1:uniform:prime-budget} gives
\[
 \log A(3N;\mathcal P)
 \le H-\frac{Hh}{\log Y}
                       +O_{K,C}(H/\log Y).
\]
Since \(\log Y=3h-\log h+O_{K,C}(1)\), the
right side equals the bound in
\eqref{p1:uniform:branching-rate}.

Exclude centres whose windows overlap an old window,
or contain a value whose square class is supported
on \(\mathcal P\). This costs at most
\(j(2L+1)+BA(3N;\mathcal P)\), with the additional
polylogarithmic factors absorbed in the error term.
Every value in a remaining window has an odd prime
outside \(\mathcal P\). Such a prime exceeds the
diameter and occurs nowhere else in the window.
These private coordinates form an identity minor
in the new valuation matrix and are independent
of all old values. This proves the exact
independence assertion.
\end{proof}

\begin{corollary}[A single pair at all critical orders]
\label{p1:uniform:one-pair}
Let \(A_N=2\sum_{e\in E_N}q_e\).
Construct ordered \(r\)-tuples by choosing the
two positions of a good pair and then adding
windows in position order, always outside the
exceptional set of
Proposition~\ref{p1:uniform:branching}.
Their total positive mass \(B_{r,N}\) satisfies,
uniformly for \(2\le r\le m\le KH/h\),
\[
 B_{r,N}=\binom r2\lambda_N^{r-2}A_N
                       \{1+O_{K,C}(rD_N^*/N)\}.
\]
Furthermore,
\[
 \sum_{r=2}^{m}\frac{2^{r-1}}{r!}
   \left|B_{r,N}-\binom r2\mu_N^{r-2}\alpha_N\right|
       =o_C(a_N).
\]
\end{corollary}

\begin{proof}
For a fixed pair each added window contributes
exactly \(p\). The number of extensions is
between \((N-D_N^*)^{r-2}\) and \(N^{r-2}\);
\(mD_N^*/N\to0\). The pair is the unique
dependent pair in the resulting tuple, so its
choice of positions is unique. This proves the
first formula.
The weighted branching error sums to
\[
 O_{K,C}\left(A_N\frac{D_N^*}{N}
       \sum_{r\ge2}\frac{2^{r-1}}{r!}
                r\binom r2\Lambda^{r-2}\right)=o(a_N),
\]
where \(\Lambda\) bounds the means.
The pair and mean estimates replace \(A_N\)
by \(\alpha_N\), and \(\lambda_N\) by \(\mu_N\),
at the same aggregate precision.
\end{proof}

This corollary controls a selected positive family.
Its independence guarantee is uniform at growing
orders, but it does not by itself estimate the full
factorial moments at those orders.

\subsection{Exact moments of the modified count}

Let \(E'_N\) be the event in Lemma~\ref{p1:core:exceptional},
let \(A=\{x\in G:J_x=1\}\), and put
\(\widetilde Z=Z_G\ind_{(E'_N)^c}\).
For \(r\ge1\), define
\[
 D_r(A)=\#\{I\subseteq A:|I|=r,\ I\text{ direct}\},
 \qquad L'_r=\E[D_r(A);E'_N].
\]
For \(r\ge2\), also put
\[
 P_r(A)=\sum_{e\in E_N}J_e
       \#\{U\subseteq A\setminus e:U\in\Gamma_e(r-2)\},
 \qquad H_r=\E[P_r(A);E'_N],
\]
and set \(H_1=0\).
Here \(\Gamma_e\), the independent deficits \(F_r,G_r\),
and \(\beta_N=\sum_eq_e\) are those used in the
stopped expansion.

\begin{proposition}[Moment decomposition and pair leakage]
\label{p1:uniform:modified-moments}
For every integer \(r\ge1\),
\begin{equation}
 \E\binom{\widetilde Z}{r}
   =\binom nrp^r-F_r-L'_r
     +\beta_N\binom{n-2}{r-2}p^{r-2}-G_r-H_r,
 \label{p1:uniform:modified-identity}
\end{equation}
where the last three terms are zero for \(r=1\).
All four deficits \(F_r,G_r,L'_r,H_r\) are
nonnegative. Moreover,
\[
 \sum_{r\ge2}2^rG_r
       \ll_\varepsilon N^{-1+\varepsilon},
 \qquad
 \sum_{r\ge2}2^rH_r=o(a_N).
\]
\end{proposition}

\begin{proof}
Outside \(E'_N\), an active family is either direct
or contains the unique active good pair and is
direct after its compression. Thus
\[
 \binom{Z_G}{r}=D_r(A)+P_r(A)
              \qquad\text{on }(E'_N)^c.
\]
The expectations without restriction to this event
are, respectively,
\[
 \E D_r(A)=\binom nrp^r-F_r,\qquad
 \E P_r(A)
     =\beta_N\binom{n-2}{r-2}p^{r-2}-G_r.
\]
Subtract the two contributions on \(E'_N\) to
obtain \eqref{p1:uniform:modified-identity}.
The bound for \(G_r\) was proved in
Lemma~\ref{p1:close:deficit-bounds}.

For each fixed \(r\), \(P_r(A)\le C_rZ_G^r\).
H\"older's inequality and the fixed moments give
\[
 H_r\le C_r(\E Z_G^{8r})^{1/8}
                  \P(E'_N)^{7/8}
       \ll_{r,\varepsilon}N^{-7/10+7\varepsilon/8}
       =o(a_N)
\]
when the margin is sufficiently small.
To sum all \(r\), use directness before taking
expectations:
\[
 \sum_{r\ge2}4^rH_r
 \le\sum_{r\ge2}4^r\E P_r(A)
 \le16\beta_N(1+4p)^{n-2}
       \ll_C a_N.
\]
For any fixed \(K\), the terms \(2\le r\le K\)
are \(o(a_N)\), and the remaining \(2^r\)-weighted
sum is \(O_C(2^{-K}a_N)\).
Take first \(N\to\infty\), then \(K\to\infty\).
This proves the asserted small \(o\), without
giving it an additional polynomial rate.
\end{proof}

The direct leakage \(L'_r\) and the deterministic
deficit \(F_r\) are different quantities.
The small probability of \(E'_N\) gives a bound
for \(L'_r\) at each fixed order, but does not
on its own control their exponentially weighted
sum. The stopped expansion avoids paying that
whole sum: it keeps the suffix inside a
positive kernel at the first excluded prefix.

\subsection{Replicas with common small-prime signs}

Fix a threshold \(Y\). Keep the signs at primes
\(\ell\le Y\) common, and give each member of a
partition \(\pi\) an independent copy of the signs
at primes \(\ell>Y\). Occurrences may have repeated
centres; their formal coordinates remain distinct.
For the resulting valuation map let \(K_\pi\)
be its kernel on all formal occurrence coordinates, before quotienting
by the internal relations of individual windows. Thus internal local
relations are included in \(K_\pi\). Let
\(\eta_\pi\in\{0,1\}\) denote affine compatibility.
If instead each local block is first replaced by its effective image,
the relative-kernel formula uses the product of the actual marginals
\(\prod_i p_{x_i}\) in place of the nominal factor \(p^r\).
Repeated centres retain distinct formal occurrence coordinates.

\begin{proposition}[The replica formula and positive propagation]
\label{p1:uniform:replicas}
For \(r\) prescribed occurrences,
\begin{equation}
 \E\prod_{A\in\pi}
     \P(J_{x_i}=1\ (i\in A)\mid\mathcal F_Y)
       =p^r\eta_\pi2^{\dim K_\pi}.
 \label{p1:uniform:replica-formula}
\end{equation}
After any fixed number of compatible colored
sources, the sum of the conditional probabilities
of targets is \(O_\varepsilon(N^\varepsilon)\),
uniformly in \(Y\) and the coloring.
Here conditioning means prescribing these source events
in the replicated space; the common small-prime signs
are still averaged. Thus a positive random sum
\[
 R=\sum_{\boldsymbol x\in\mathcal A}w_{\boldsymbol x}
       \prod_{A\in\pi}
          \P(J_{x_i}=1\ (i\in A)\mid\mathcal F_Y),
\]
with fixed tuple length, fixed partition \(\pi\), and
nonnegative deterministic weights bounded by a fixed
power of \(\log N\),
satisfies for every fixed integer \(q\ge1\)
\begin{equation}
 \E R^q\ll_\varepsilon N^\varepsilon\E R,
 \label{p1:uniform:replica-moment}
\end{equation}
when its index set is a subfamily of all ordered
tuples of the stated fixed length.
\end{proposition}

\begin{proof}
Conditionally on the common small-prime signs,
the independent large-prime copies realize the
product of conditional probabilities on the left
of \eqref{p1:uniform:replica-formula}.
Fourier inversion in this larger binary space
gives its right side, with its compatibility
factor retained.

Merging colors adds prime-parity equations
together, so a relation in the colored system
is also a relation in the merged one. The
projection of its kernel to the target block
is contained in the corresponding merged
projection. Consequently the target's colored
rank loss is no larger than its uncolored rank
loss. Proposition~\ref{p1:fixed:targets} and the
proof of \eqref{p1:fixed:conditional} apply.
Targets whose centre repeats a source account
for only a fixed number of choices and have
probability at most one.

Expand \(R^q\) positively, retain the first
replica product as a seed, and sum all further
occurrences successively. Each step has the
conditional bound just proved. Repeated centres
are treated as above. There are only a fixed
number of steps; splitting \(\varepsilon\)
among them and absorbing the fixed powers of
\(\log N\) proves \eqref{p1:uniform:replica-moment}.
\end{proof}

\paragraph{A local-kernel example.}
For a single start at \(x=18\) of length \(L=10\), the ten affine
rows have rank nine. At prime cutoff \(28\), the unique compatible
assignment makes prime \(17\) negative and all other prime signs
positive, so \(p_{18}=2^{-9}\). The formal relation kernel has
dimension one, giving \(2^{-10}2^1=2^{-9}\) in the nominal formula.
The effective relative nullity of a singleton is zero; its formula
therefore retains \(p_{18}\). The internal even square relation is
visible in \(18\cdot24\cdot25\cdot27=540^2\).
This is a finite algebraic illustration, not an assertion that these
parameters lie beyond an asymptotic threshold.

The same reasoning permits measurable gates in
positive estimates. For example, suppose
\(\E R\ll_\delta N^{-\eta+\delta}\) for every
\(\delta>0\), and let \(\mathcal A\in\mathcal F_Y\).
For an integer \(k\ge1\), interpolation gives
\[
 \E R^{1+1/k}\le(\E R)^{1-1/k}(\E R^2)^{1/k}
       \ll_\delta N^\delta\E R.
\]
H\"older's inequality with conjugate exponents
\(1+1/k\) and \(k+1\), followed by the fixed
moment of order \(d(k+1)\), gives
\(\E[RZ_G^d]\ll_\delta
 N^{-\eta k/(k+1)+O(\delta)}\).
Choose \(k\) after the desired error margin.
Since the gate is at most one, for fixed \(d\) this yields
\[
 \E[R\,Z_G^d\,\ind_{\mathcal A}]
       \ll_\varepsilon N^{-\eta+\varepsilon}.
\]
Here \eqref{p1:uniform:replica-moment} supplies the
second moment of \(R\). The estimate is an
unconditional weighted expectation; it makes
no division by \(\P(\mathcal A)\).
